\chapter*{Resumen}
\addcontentsline{toc}{chapter}{Resumen}
\markboth{Resumen}{Resumen}

Este Trabajo de Fin de Grado plantea una comparación experimental entre dos estrategias para analizar electrocardiogramas representados como imagen cuando los datos etiquetados son escasos. La pregunta central es directa: con un mismo presupuesto de \(K\) ejemplos etiquetados, ¿cuándo compensa entrenar un modelo de visión específico (CNN) y cuándo resulta preferible aprovechar el conocimiento multimodal de un modelo de visión y lenguaje (VLM) mediante aprendizaje en contexto? El caso clínico motivador es la sospecha de síndrome de Brugada, una canalopatía hereditaria del canal de sodio asociada a muerte súbita cardíaca.

El trabajo se enmarca estrictamente como apoyo al triaje clínico. En la práctica asistencial, cuando un ECG sugiere un patrón sospechoso de Brugada, se inicia un circuito que incluye lectura especializada del trazado, valoración de síntomas y antecedentes familiares, posible recolocación de las derivaciones V1 y V2 en espacios intercostales superiores y, en pacientes seleccionados con sospecha persistente y ECG basal no concluyente, una prueba farmacológica con bloqueantes del canal de sodio (ajmalina, flecainida, procainamida o pilsicainida) en un entorno monitorizado con capacidad de reanimación. Esta prueba puede desenmascarar un patrón tipo 1, pero debe interpretarse dentro del contexto clínico y tras excluir fenocopias. El modelo desarrollado no diagnostica ni descarta pacientes: prioriza trazados que merecen revisión temprana por especialistas.

La memoria integra dos dominios. El primero es un simulador controlado de morfología QRS/ST que genera imágenes de V1, V2 y V3 con etiquetas derivadas de reglas observables. El segundo es el conjunto Brugada HUCA, procesado desde registros clínicos reales con 317 pacientes incluidos, 951 imágenes y una etiqueta binaria de caso Brugada confirmado frente a control. La evaluación se organiza por paciente mediante validación leave-one-out, presupuestos reducidos de \(K=\{2,4,8,16,32\}\) ejemplos y tres semillas por presupuesto.

La rama CNN cumple una función deliberada: demostrar empíricamente que el enfoque supervisado clásico no funciona con tan pocos datos, ni siquiera con técnicas de adaptación de dominio. ResNet18 aprende de forma clara en el simulador, con exactitud equilibrada de 0,848 para \(K=32\), sensibilidad 0,883 y especificidad 0,812. Sin embargo, en Brugada HUCA el rendimiento cae cerca del azar: la mejor exactitud equilibrada del modelo base es 0,516 con \(K=16\). La adaptación de dominio (SSL, CORAL, MMD, DANN) mejora el mejor punto hasta 0,549 con MMD y \(K=32\), una mejora insuficiente para considerar el sistema robusto.

Esta evidencia negativa constituye una contribución central del TFG. No se trata de un fallo del proyecto, sino del resultado que fundamenta la hipótesis principal: si entrenar un modelo visual específico no basta con pocos datos, el conocimiento multimodal previo de un VLM podría compensar esa escasez. Sobre esa base, la rama VLM/ICL se desarrolla como comparador principal: Gemma 4 y MedGemma reciben las mismas imágenes y los mismos presupuestos \(K\), pero los ejemplos se incluyen como demostraciones en el prompt sin actualizar pesos. El documento incorpora el protocolo completo, los mensajes, la validación JSON, los controles y las tablas de resultados preparadas; solo las celdas numéricas VLM quedan marcadas como \emph{por ejecutar} para no inventar inferencias.
